\documentclass[10pt,a4paper]{amsart}
\usepackage[margin=19mm]{geometry}
\usepackage{amsmath,amssymb,mathtools}
\usepackage[scr=rsfs,cal=euler]{mathalfa}
\usepackage{xfrac}
\usepackage[unicode,hidelinks,bookmarks=false]{hyperref}
\newcommand{\ie}{i.e.\ }
\newcommand{\cf}{cf.\ }
\newcommand{\resp}{respectively\ }
\newcommand{\Subsection}{\subsection}
\providecommand{\nobreakdash}{\nobreak\hbox{-}}
\setlength{\emergencystretch}{2em}
\multlinegap0pt
\usepackage{bm}
\usepackage{bbm}
\usepackage{dsfont}
\usepackage{xspace}
\usepackage{cancel}
\usepackage{mathtools}
\usepackage{amsthm}
\makeatletter
\@ifundefined{theorem}{\newtheorem{theorem}{Theorem}}{}
\@ifundefined{lemma}{\newtheorem{lemma}[theorem]{Lemma}}{}
\@ifundefined{proposition}{\newtheorem{proposition}[theorem]{Proposition}}{}
\makeatother
\DeclareMathOperator{\ord}{ord}
\newcommand{\C}{\mathbf C}
\newcommand{\G}{\Gamma_0(3)}
\newcommand{\Up}{\Lambda_p}
\newcommand{\Vp}{V_p}
\newcommand{\Z}{\mathbf Z}
\title[A full $p^{4r}$ supercongruence tower for a level-three symm]{A full $p^{4r}$ supercongruence tower for a level-three symmetric-cube hypergeometric sequence}
\author{Alex Shvets}
\date{Source: arXiv:2606.15462v1 (CC BY 4.0)}
\begin{document}
\maketitle

\noindent\emph{Source and licence.} A full $p^{4r}$ supercongruence tower for a level-three symmetric-cube hypergeometric sequence. Version arXiv:2606.15462v1 (CC BY 4.0), distributed under the Creative Commons Attribution 4.0 International licence, as recorded in the arXiv licence metadata for this exact version and shown on the abstract page https://arxiv.org/abs/2606.15462v1. Full rights notice: licenses/arxiv-2606.15462-CC-BY-4.0.txt.\par\smallskip

\noindent\textit{Editorial scope.} Counted unit: the Theorem whose source label is \texttt{thm:sparse-tower} in the article (source0.tex lines 848--853, source label \texttt{thm:sparse-tower}), delivered together with the article's own complete original proof of it (source0.tex lines 855--949, one \texttt{proof} environment of 95 lines). No mathematics is written, repaired or re-derived: statement and proof are the article's own text line for line. Required context retained uncounted: prop:fricke-C (lines 450--473); prop:polelowering (lines 563--570); prop:hecke-prime-power (lines 617--627); prop:fricke-hecke (lines 666--672); lem:order-estimate (lines 733--742); lem:sparse-pole (lines 785--791); lem:G-F-comparison (lines 804--814); lem:G-cusps (lines 838--840). Every definition, display and label of those lines is retained unchanged. Omitted from this package and not redistributed: 1--449, 474--562, 571--616, 628--665, 673--732, 743--784, 792--803, 815--837, 841--847, 854--854, 950--1024. Nothing inside a retained excerpt is omitted or altered: every retained body is delivered exactly as the article's lines are written, so no omission receipt of any kind accompanies this package. Citations: the retained excerpts carry no citation command at all, so no bibliography entry of the article is transcribed and no reference is redistributed. Reference adaptation: none was needed and none was made; every reference inside the delivered unit resolves inside it, so no unresolved reference or citation is produced. Printed slips kept literal: no printed word, symbol, hypothesis or number of the counted statement or of its proof was altered. Layout and class: the article's own class is \texttt{amsart} with packages including amsmath, amssymb, amsthm, mathtools, fontenc, inputenc, lmodern, microtype, hyperref; the wrapper is the lane's pinned \texttt{amsart} editorial wrapper at 10pt on A4, which restates the theorem-like environments the retained text needs and adds the article's own macro definitions that the retained text needs (\texttt{\textbackslash C}, \texttt{\textbackslash G}, \texttt{\textbackslash Up}, \texttt{\textbackslash Vp}, \texttt{\textbackslash Z}, \texttt{\textbackslash ord}), with extra wrapper packages (bm, bbm, dsfont, xspace, cancel, mathtools). The delivered numbering is the unit's own. Assets: no retained span contains a figure environment, a graphics-inclusion command, a PDF-page inclusion, a table or a TikZ picture; no image file is copied into this package, the package declares no asset dependency and no image file is redistributed with it. Licence: version arXiv:2606.15462v1 of this article is distributed under the Creative Commons Attribution 4.0 International licence, as recorded in the arXiv licence metadata for this exact version and shown on the abstract page https://arxiv.org/abs/2606.15462v1. A full rights notice is delivered at licenses/arxiv-2606.15462-CC-BY-4.0.txt. Characters: every retained excerpt is pure ASCII, so the delivered engine, XeLaTeX with its default Latin Modern text font, reports no missing character.
\begin{proposition}[Fricke identities]\label{prop:fricke-C}
One has
\[
 t|_0W_3=\frac{3^{-6}}{t},
 \qquad
 \Theta|_3W_3=27i\,t\Theta,
 \qquad
 \left(\frac{Dt}{t}\right)\!\Big|_2W_3=-\frac{Dt}{t}.
\]
Consequently
\[
 C|W_3=-27i\,tC,
 \qquad
 (tC)|W_3=-\frac{i}{27}C.
\]
Moreover
\[
 \ord_0(C)=1,
 \qquad
 \ord_0(tC)=0,
 \qquad
 \ord_0\left(\frac{C}{t^j}\right)=j+1\quad(j\ge0).
\]
\end{proposition}

\begin{proposition}[Pole-lowering expansion]\label{prop:polelowering}
Let $f\in M_5^!(\G,\chi_3)$ have poles only at the cusp $\infty$, and suppose the pole order at $\infty$ is at most $N$.  Then
\[
 f=\beta_{-1}tC+\sum_{j=0}^{N}\alpha_j\frac{C}{t^j}
\]
for some $\beta_{-1},\alpha_0,\ldots,\alpha_N\in\C$.
If the $q$-expansion of $f$ is in $\Z_{(p)}((q))$, then all these coefficients lie in $\Z_{(p)}$.
\end{proposition}

\begin{proposition}[Prime-power Hecke expansion]\label{prop:hecke-prime-power}
Let $p\ge5$, put $\chi=\chi_3(p)$, and let $a\ge0$.  On $q$-expansions of weakly holomorphic forms of weight $5$ and character $\chi_3$ on $\G$,
\[
 T_{p^a}=
 \sum_{i=0}^{a}\chi^i p^{4i}\Vp^i\Up^{a-i}.
\]
In particular,
\[
 T_p=\Up+\chi p^4\Vp.
\]
\end{proposition}

\begin{proposition}[Fricke--Hecke intertwining]\label{prop:fricke-hecke}
For every $a\ge0$ and every prime $p\ge5$,
\[
 T_{p^a}W_3=\chi_3(p)^a W_3T_{p^a}
\]
on $M_5^!(\G,\chi_3)$.
\end{proposition}

\begin{lemma}[Order estimate]\label{lem:order-estimate}
Let $a\ge0$, $p\ge5$, and let
\[
 g(q)=\sum_{n\ge N}a_nq^n
\]
with $N\ge0$.  Then
\[
 \ord_\infty(T_{p^a}g)\ge\left\lceil\frac{N}{p^a}\right\rceil.
\]
\end{lemma}

\begin{lemma}[Pole order of the sparse defect]\label{lem:sparse-pole}
For every $s\ge0$ and $m\ge1$,
\[
 \mathcal F_{m,s}=O(q^{-m+1})
\]
at $\infty$.
\end{lemma}

\begin{lemma}[Comparison between Hecke and sparse defects]\label{lem:G-F-comparison}
Assume that for all $t<s$ and all $M\ge1$ one has
\[
 \mathcal F_{M,t}\equiv0\pmod {p^{4(t+1)}}.
\]
Then, for every $m\ge1$,
\[
 \mathcal G_{m,s}\equiv \mathcal F_{m,s}\pmod {p^{4(s+1)}}.
\]
For $s=0$ the same conclusion holds without any hypothesis.
\end{lemma}

\begin{lemma}[Modularity and cusps of $\mathcal G_{m,s}$]\label{lem:G-cusps}
For all $m\ge1$ and $s\ge0$, the form $\mathcal G_{m,s}$ belongs to $M_5^!(\G,\chi_3)$ and has poles only at the cusp $\infty$.
\end{lemma}

\begin{theorem}[Sparse Cartier tower]\label{thm:sparse-tower}
For every prime $p\ge5$, every $m\ge1$, and every $s\ge0$,
\[
 \mathcal F_{m,s}\equiv0\pmod {p^{4(s+1)}}.
\]
\end{theorem}

\begin{proof}
We prove the statement by induction on $s$, uniformly in $m$.
Assume first that the statement is known for all lower values $t<s$; when $s=0$ this assumption is empty.
By Lemma~\ref{lem:G-F-comparison},
\[
 \mathcal G_{m,s}\equiv\mathcal F_{m,s}\pmod {p^{4(s+1)}}.
\]
By Lemma~\ref{lem:sparse-pole},
\[
 \mathcal F_{m,s}=O(q^{-m+1}).
\]
Therefore
\[
 [q^{-j}]\mathcal G_{m,s}\equiv0\pmod {p^{4(s+1)}}
 \qquad(j\ge m).
\]

Moreover, $\mathcal G_{m,s}$ has $q$-expansion in $\Z_{(p)}((q))$: this follows from $f_N\in q^{-N}\Z[[q]]$ and the prime-power expansion of Proposition~\ref{prop:hecke-prime-power}.  Therefore, by Lemma~\ref{lem:G-cusps} and Proposition~\ref{prop:polelowering}, there is a finite expansion
\[
 \mathcal G_{m,s}=\beta_{-1}tC+\sum_{j=0}^{J}\alpha_j\frac{C}{t^j}
\]
with coefficients in $\Z_{(p)}$.  The triangular form of the basis at $\infty$ gives, by descending induction on $j$,
\[
 \alpha_j\equiv0\pmod {p^{4(s+1)}}
 \qquad(j\ge m).
\]
Indeed, the coefficient of $q^{-j}$ in $C/t^j$ is $1$, and only the terms $C/t^\ell$ with $\ell\ge j$ can contribute to $q^{-j}$.

It remains to kill the coefficients with $j<m$ and the coefficient of $tC$.  Apply $W_3$.  By Proposition~\ref{prop:fricke-hecke},
\[
 \mathcal G_{m,s}|W_3
 =\chi_3(p)^{s+1}T_{p^{s+1}}(f_{mp^{s+1}}|W_3)
 -\chi_3(p)^sT_{p^s}(f_{mp^s}|W_3).
\]
By Proposition~\ref{prop:fricke-C},
\[
 f_N|W_3=\left(\frac{C}{t^N}\right)\Big|W_3
 =-27i\,3^{6N}Ct^{N+1}.
\]
Thus
\[
 \ord_\infty(f_{mp^{s+1}}|W_3)=mp^{s+1}+1,
 \qquad
 \ord_\infty(f_{mp^s}|W_3)=mp^s+1.
\]
Lemma~\ref{lem:order-estimate} gives
\[
 \ord_\infty T_{p^{s+1}}(f_{mp^{s+1}}|W_3)
 \ge\left\lceil\frac{mp^{s+1}+1}{p^{s+1}}\right\rceil=m+1
\]
and
\[
 \ord_\infty T_{p^s}(f_{mp^s}|W_3)
 \ge\left\lceil\frac{mp^s+1}{p^s}\right\rceil=m+1.
\]
Therefore
\[
 \ord_\infty(\mathcal G_{m,s}|W_3)\ge m+1.
\]

On the other hand,
\[
 (tC)|W_3=-\frac{i}{27}C=-27i\,3^{-6}C
\]
and, for $j\ge0$,
\[
 \left(\frac{C}{t^j}\right)\Big|W_3
 =-27i\,3^{6j}Ct^{j+1}.
\]
Hence
\[
 \mathcal G_{m,s}|W_3
 =-27i\left(\beta_{-1}3^{-6}C+
 \sum_{j=0}^{J}\alpha_j3^{6j}Ct^{j+1}\right).
\]
The forms
\[
 C,\ Ct,\ Ct^2,\ldots
\]
have distinct orders $0,1,2,\ldots$ at $\infty$.  Since $\ord_\infty(\mathcal G_{m,s}|W_3)\ge m+1$, the coefficients of $C,Ct,\ldots,Ct^m$ vanish exactly.  Therefore
\[
 \beta_{-1}=0,
 \qquad
 \alpha_j=0\quad(0\le j\le m-1).
\]
Combining these exact equalities with the congruences for $j\ge m$ gives
\[
 \mathcal G_{m,s}\equiv0\pmod {p^{4(s+1)}}.
\]
Since $\mathcal G_{m,s}\equiv\mathcal F_{m,s}\pmod {p^{4(s+1)}}$, we conclude
\[
 \mathcal F_{m,s}\equiv0\pmod {p^{4(s+1)}}.
\]
This completes the induction on $s$.
\end{proof}

\end{document}
